\documentclass[UTF8,fontset=fandol,11pt]{ctexart}
\usepackage[a4paper,margin=25mm]{geometry}
\usepackage{amsmath,amssymb,amsthm,mathtools}
\usepackage{xcolor}
\usepackage{booktabs,array,enumitem,fancyhdr,hyperref}
\hypersetup{colorlinks=true,linkcolor=blue!50!black,urlcolor=blue!50!black}
\setlength{\emergencystretch}{2em}
\setlength{\parskip}{0.3em}
\linespread{1.16}
\allowdisplaybreaks
\newtheorem{theorem}{定理}[section]
\newtheorem{proposition}[theorem]{命题}
\newtheorem{lemma}[theorem]{引理}
\theoremstyle{definition}
\newtheorem{definition}[theorem]{定义}
\theoremstyle{remark}
\newtheorem{remark}[theorem]{说明}
\newcommand{\dd}{\,\mathrm d}
\newcommand{\ii}{\mathrm i}
\newcommand{\Rea}{\operatorname{Re}}
\newcommand{\Ima}{\operatorname{Im}}
\newcommand{\Li}{\operatorname{Li}}
\newcommand{\supp}{\operatorname{supp}}
\pagestyle{fancy}
\fancyhf{}
\fancyhead[L]{\small 学习笔记}
\fancyhead[R]{\small 初稿}
\fancyfoot[C]{\thepage}
\setlength{\headheight}{15pt}
\author{}
\date{2026年10月22日初稿}
\makeatletter
\renewcommand{\maketitle}{%
  \thispagestyle{plain}%
  \begin{center}
    {\LARGE\bfseries \@title\par}
    \vspace{0.6em}
    {\small \@date\par}
  \end{center}
  \vspace{0.6em}}
\makeatother
\newcommand{\draftnotice}{\noindent{\small
本文为供逐篇审阅的学习笔记初稿，不是新增研究成果。
文中区分直接证明与引用结果；正文不依赖本文。
本轮已作公式、条件及排版初步核查，尚待作者逐句审定。}\par\medskip}

\usepackage{pgfplots}
\usepackage{float}
\pgfplotsset{compat=1.18}
\date{2026年10月22日初稿；10月23日增补示意图}
\title{渐近级数为什么越加越不准确}
\begin{document}
\maketitle
\draftnotice
\section{先澄清两种不同的极限}
渐近展开研究的是：固定截断阶数，让变量趋向指定极限，余项比保留项小到什么程度。
收敛级数研究的是：固定变量，让项数趋于无穷，部分和是否趋向函数值。
两者交换量词便会变成不同命题。
本篇用一个有精确余项的积分说明区别，再讨论对数积分中的阶乘系数。

\section{一个可完全控制的例子}
对 $t>0$ 定义
\[
 F(t)=\int_0^\infty\frac{e^{-u}}{1+tu}\dd u.
\]
有限几何恒等式对任意正整数 $N$ 都成立：
\[
 \frac1{1+tu}=\sum_{k=0}^{N-1}(-tu)^k+\frac{(-tu)^N}{1+tu}.
\]
逐项积分的和是有限和，故不涉及无穷换序：
\[
 F(t)=\sum_{k=0}^{N-1}(-1)^kk!t^k+R_N(t),\qquad
 R_N(t)=(-t)^N\int_0^\infty\frac{e^{-u}u^N}{1+tu}\dd u.
\]
因此
\[
 0<(-1)^NR_N(t)<N!t^N.
\]
固定 $N$，令 $t\to0^+$，就得到各有限阶的渐近控制。
然而固定任何 $t>0$ 时，相邻项的绝对值之比为 $(k+1)t\to\infty$，
项本身最终不趋于零，所以无穷级数发散。

\section{项最小时与误差界最小时}
界 $b_N=N!t^N$ 满足 $b_{N+1}/b_N=(N+1)t$。
当 $N$ 小于约 $1/t$ 时界逐步减小，超过这一范围后界变大。
这解释了为什么适当增加项数可能有益，而无限增加项数并非合理策略。

但须区分：$b_N$ 是余项的上界，其最小点不必恰好等于实际误差 $|R_N|$ 的最小点。
由积分的精确表达式，可以作进一步比较；不应把“截在最小项附近”说成无需条件的普遍定理。
本例每个 $R_N$ 的符号已知，偶数项数和奇数项数还分别提供一侧界。

\subsection*{固定参数下的截断对照}
\newcommand{\trialF}{0.9192568286758049}
\newcommand{\trialBoundN}{10}
\newcommand{\trialErrorN}{11}
\newcommand{\trialBoundMin}{0.000217269660510822937}
\newcommand{\trialErrorMin}{0.000108559235819575947}

图~\ref{fig:asymptotic-truncation} 固定 $t=0.095$，比较保留 $N$ 项后的
上界 $b_N=N!t^N$ 与实际误差
$E_N=|F(t)-\sum_{k=0}^{N-1}(-1)^kk!t^k|$。
本次计算得 $F(0.095)\approx\trialF$。
在图示的 $1\le N\le25$ 内，上界在 $N=\trialBoundN$ 处最小，
实际误差在 $N=\trialErrorN$ 处最小；两者不应混称为同一个最优截断。
\begin{figure}[H]
\centering
\begin{tikzpicture}
\begin{semilogyaxis}[
  width=14.6cm,height=7cm,xmin=1,xmax=25,ymin=0.00001,ymax=1,
  xlabel={保留项数 $N$},ylabel={绝对误差或上界},
  xtick={1,5,10,15,20,25},grid=major,grid style={gray!20},
  tick label style={font=\small},label style={font=\small},
  legend style={font=\small,draw=none,at={(0.02,0.98)},anchor=north west},
  axis lines=left]
\addplot[black,thick,dashed,mark=square*,mark size=1.6pt]
  table[x=N,y=bound,col sep=comma] {figures/truncation-data.csv};
\addlegendentry{$b_N=N!t^N$}
\addplot[blue!65!black,thick,mark=*,mark size=1.6pt]
  table[x=N,y=error,col sep=comma] {figures/truncation-data.csv};
\addlegendentry{$E_N=|R_N(t)|$}
\draw[gray!80,densely dotted] (axis cs:\trialBoundN,\trialBoundMin)
 -- (axis cs:\trialBoundN,0.0008);
\draw[blue!65!black,densely dashed] (axis cs:\trialErrorN,\trialErrorMin)
 -- (axis cs:\trialErrorN,0.00003);
\node[anchor=south,font=\footnotesize,fill=white,inner sep=2pt]
  at (axis cs:\trialBoundN,0.0008) {$b_N$ 最小：$N=\trialBoundN$};
\node[anchor=north west,font=\footnotesize,fill=white,inner sep=2pt]
  at (axis cs:\trialErrorN,0.00003) {$E_N$ 最小：$N=\trialErrorN$};
\end{semilogyaxis}
\end{tikzpicture}

\caption{固定 $t=0.095$ 时的误差上界与实际误差，对数纵轴。
每个点对应一个整数项数 $N$，连线只帮助比较趋势；竖虚线标出各自的最小点。
数值由精确余积分计算，并与部分和之差交叉复核，不是区间认证。
本图只对应 $F(t)$，不能将其误差界直接用于后文的 $\Li$。}
\label{fig:asymptotic-truncation}
\end{figure}
两条曲线先降后升，说明在固定参数下增加项数可能先改善、后恶化近似；
前节关于每个固定 $N$、$t\to0^+$ 的结论则仍然成立。
图示仅比较这一组参数，并不把最小项附近截断提升为一般定理。

\section{对数积分的有限阶展开}
在 $x>2$ 上定义 $\Li(x)=\int_2^x\dd v/\log v$。
记 $L=\log x$。分部积分得到
\[
 \Li(x)=\frac{x}{L}-\frac2{\log2}
        +\int_2^x\frac{\dd v}{(\log v)^2}.
\]
递推使用
\[
 \int_2^x\frac{\dd v}{(\log v)^j}
 =\frac{x}{(\log x)^j}-\frac2{(\log2)^j}
  +j\int_2^x\frac{\dd v}{(\log v)^{j+1}},
\]
便得对每个固定 $N\ge1$ 的精确恒等式
\[
 \Li(x)=x\sum_{k=0}^{N-1}\frac{k!}{L^{k+1}}
       -C_N+N!\int_2^x\frac{\dd v}{(\log v)^{N+1}},
 \qquad
 C_N=2\sum_{k=0}^{N-1}\frac{k!}{(\log2)^{k+1}}.   \tag{1}
\]
下限常数有明确表达式，不能在讨论固定 $x$ 的数值时擅自丢掉。

\section{余项的阶为何成立}
对固定整数 $j\ge1$，令 $I_j(x)=\int_2^x(\log v)^{-j}\dd v$。
$I_j(x)$ 与 $x/(\log x)^j$ 均趋于无穷，且后者在 $x>e^j$ 时导数为
\[
 \frac1{(\log x)^j}\left(1-\frac j{\log x}\right)>0.
\]
由 l'Hôpital 法则，
\[
 I_j(x)\sim \frac{x}{(\log x)^j}.
\]
取 $j=N+1$ 并在式(1)中固定 $N$，可得
\[
 \Li(x)=x\sum_{k=0}^{N-1}\frac{k!}{(\log x)^{k+1}}
       +O_N\left(\frac{x}{(\log x)^{N+1}}\right).
\]
这里 $C_N$ 被吸收到余项，是因为 $N$ 固定且 $x/(\log x)^{N+1}\to\infty$。
隐含常数依赖 $N$；不能据此允许 $N$ 任意随 $x$ 增长。

\section{如何避免对发散展开的误用}
固定 $x$ 时，形式项之比为 $(k+1)/\log x$，最终增大；
$\Li$ 的形式无穷展开不收敛。它也不是上一个积分 $F(t)$ 的交错级数，
因而不能直接套用交错余项界。若要优化 $\Li$ 的数值截断，
应另用具有相应参数范围的余项估计，或直接计算式(1)的余积分。

渐近展开的价值并未因发散而消失：任何固定阶都提供可证明的大变量近似。
需要舍弃的是“项数越多必越准”的推断，而不是有限阶渐近本身。

\section*{参考依据与后续核验}
本文的两个余项表达式均由有限恒等式或有限次分部积分得到。
相关背景可参看本书《素数分布：从古典经验公式到解析数论》及 Stirling 展开的讨论。
绘图数据、参数和生成程序随本文提供。数据由 \texttt{mpmath} 在80位工作精度下计算，
并以120位精度复算；截断误差同时用精确余积分与部分和之差核对。
后续增加数值例子时仍须报告变量、项数及实际误差，不能只画部分和而不比较函数值。
\end{document}
